\documentclass[11pt]{article}

\usepackage[T1]{fontenc}
\usepackage{amsmath,amssymb,amsthm}
\usepackage{hyperref}

\title{Fixed-Misspecification RMSEA Bias:\\A Curvature Check}
\author{}
\date{\today}

\newcommand{\E}{\mathbb{E}}
\newcommand{\Var}{\operatorname{Var}}
\newcommand{\tr}{\operatorname{tr}}
\newcommand{\argmin}{\operatorname*{arg\,min}}
\newcommand{\dto}{\xrightarrow{d}}
\newcommand{\pto}{\xrightarrow{p}}

\theoremstyle{plain}
\newtheorem{theorem}{Theorem}
\newtheorem{corollary}{Corollary}
\theoremstyle{remark}
\newtheorem*{remark}{Remark}

\begin{document}
\maketitle

\section*{Purpose}

This note checks the fixed-misspecification version of the RMSEA bias
calculation. The local-drift story gives the familiar chain
\[
  df \quad\longrightarrow\quad \tr(U\Gamma),
\]
where $U$ is the first-order residual projector and $\Gamma$ is the asymptotic
covariance of the saturated moments. Under fixed model misspecification this is
not the full $n^{-1}$ bias of the fitted discrepancy. The population residual is
then $O(1)$, so two terms that are invisible under exact fit or Pitman drift move
to leading bias order:
\[
  \text{model curvature against the residual,}
  \qquad
  \text{data-dependent weight influence.}
\]
There is also a mundane but important bookkeeping term: if the first-stage
moment estimator has an $O(n^{-1})$ mean bias, that bias contributes at the same
order because the population residual is nonzero.

The result below is therefore the fixed-misspecification correction
\[
  \E[\hat F]
  =
  F_0
  +
  \frac{1}{n}
  \left\{
    \tr(\widetilde U\Gamma)
    + \beta_W
    + 2 e'W b
  \right\}
  + o(n^{-1}),
\]
with $\widetilde U$ built from the observed population Hessian, not just the
Gauss--Newton Hessian. For a fixed weight and an unbiased first-stage covariance
convention this reduces to $\tr(\widetilde U\Gamma)$.

\section{Setup}

Let $m$ be the saturated moment vector used by the estimator. For ordinary
complete-data covariance SEM, $m=\operatorname{vech}(\Sigma)$; for other
settings it may include means, thresholds, polychorics, or group blocks. Assume
the sample version has the expansion
\begin{equation}\label{eq:mexp}
  \hat m
  =
  m_0 + n^{-1/2} z + n^{-1} b_n + o_p(n^{-1}),
  \qquad
  z\dto N(0,\Gamma),
  \qquad
  \E[b_n]\to b .
\end{equation}
The vector $b$ is zero for an exactly unbiased first-stage convention, but it is
not automatically zero for every covariance divisor or saturated statistic.

Let the model be a smooth map $\theta\mapsto\mu(\theta)$ and define the
quadratic discrepancy
\begin{equation}\label{eq:Fhat}
  \hat F
  =
  \min_\theta
  r_n(\theta)'\hat W r_n(\theta),
  \qquad
  r_n(\theta)=\hat m-\mu(\theta).
\end{equation}
The pseudo-true point for the limiting weight $W$ is
\begin{equation}\label{eq:pseudotrue}
  \theta_\ast
  =
  \argmin_\theta
  \{m_0-\mu(\theta)\}'W\{m_0-\mu(\theta)\},
  \qquad
  e=m_0-\mu(\theta_\ast),
  \qquad
  F_0=e'We .
\end{equation}
Write $D=\partial\mu/\partial\theta'$ at $\theta_\ast$. The population first
order condition is
\begin{equation}\label{eq:foc}
  D'We=0 .
\end{equation}

Let $H_a$ be the Hessian matrix of the $a$th component of $\mu(\theta)$ at
$\theta_\ast$. Define
\begin{equation}\label{eq:ABK}
  A = D'WD,
  \qquad
  K_{ij} = \sum_a (We)_a (H_a)_{ij},
  \qquad
  B=A-K .
\end{equation}
Here $A$ is the Gauss--Newton part and $B$ is one half of the population Hessian
of $e'We$ with respect to $\theta$. The residual-curvature term $K$ vanishes
under correct specification and for a locally linear model. Under fixed
misspecification it is generally not negligible.

For an estimated weight, write
\begin{equation}\label{eq:Wexp}
  \hat W
  =
  W + n^{-1/2} W_1 + n^{-1} W_2 + o_p(n^{-1}),
\end{equation}
with $W_1$ the mean-zero first-order influence of the weight map. Put
\begin{equation}\label{eq:g}
  g=W_1e,
  \qquad
  P=DB^{-1}D',
  \qquad
  \widetilde U = W-WDB^{-1}D'W = W-WPW .
\end{equation}

\section{Expansion}

\begin{theorem}[Fixed-misspecification expansion]\label{thm:expansion}
Assume \eqref{eq:mexp}--\eqref{eq:Wexp}, smoothness, full column rank $D$, and
nonsingular $B$. Then
\begin{equation}\label{eq:Fexp}
  \hat F
  =
  F_0
  + n^{-1/2}\{2e'Wz + e'W_1e\}
  + n^{-1}\Phi
  + o_p(n^{-1}),
\end{equation}
where
\begin{equation}\label{eq:Phi}
  \Phi
  =
  z'\widetilde U z
  + 2g'(I-PW)z
  - g'Pg
  + e'W_2e
  + 2e'Wb_n .
\end{equation}
\end{theorem}

\begin{proof}
Let $\hat\theta-\theta_\ast=n^{-1/2}\xi+O_p(n^{-1})$. Expanding the estimating
equation $D(\hat\theta)'\hat W\{\hat m-\mu(\hat\theta)\}=0$ to first order and
using $D'We=0$ gives
\[
  (A-K)\xi = D'(Wz+W_1e),
  \qquad
  \xi=B^{-1}D'(Wz+g).
\]
Now expand the fitted residual:
\[
  r_n(\hat\theta)
  =
  e+n^{-1/2}(z-D\xi)
  +n^{-1}\{b_n-\eta-\tfrac12 H[\xi,\xi]\}
  +o_p(n^{-1}),
\]
where $\eta$ is the second-order parameter increment and $H[\xi,\xi]$ stacks
$\xi'H_a\xi$. In $r_n(\hat\theta)'\hat W r_n(\hat\theta)$, the population FOC
kills the terms involving $\eta$ and leaves
\[
  z'Wz - 2\xi'D'Wz + \xi'B\xi
  + 2g'(z-D\xi)
  + e'W_2e
  + 2e'Wb_n .
\]
Substituting $\xi=B^{-1}D'(Wz+g)$ and collecting terms gives
\eqref{eq:Phi}.
\end{proof}

\begin{corollary}[Bias]\label{cor:bias}
If $\E[W_1]=0$ and the expectations below exist, then
\begin{equation}\label{eq:bias}
  \E[\hat F]
  =
  F_0
  + \frac{1}{n}
  \left[
    \tr(\widetilde U\Gamma)
    + \beta_W
    + 2e'Wb
  \right]
  + o(n^{-1}),
\end{equation}
where
\begin{equation}\label{eq:betaW}
  \beta_W
  =
  2\E\{g'(I-PW)z\}
  - \E(g'Pg)
  + e'\E(W_2)e .
\end{equation}
\end{corollary}

For a fixed weight, $W_1=W_2=0$ and $\beta_W=0$. If additionally $b=0$, the
whole correction is just
\[
  \E[\hat F]=F_0+n^{-1}\tr(\widetilde U\Gamma)+o(n^{-1}).
\]
Under correct specification, $e=0$, so $K=0$, $\widetilde U$ reduces to the
usual first-order residual projector
\[
  U = W-WD(D'WD)^{-1}D'W ,
\]
and the bias reduces to $n^{-1}\tr(U\Gamma)$. Under normal-theory weighting
with $W=\Gamma^{-1}$, this trace is $df$.

\section{Consequences for RMSEA}

The fixed-population RMSEA target is
\[
  \varepsilon_0 = \sqrt{F_0/df}.
\]
When $F_0>0$ is fixed, the leading distribution of $\hat F$ is not noncentral
chi-square. It is ordinary root-$n$ normal, because the linear term in
\eqref{eq:Fexp} is nonzero:
\begin{equation}\label{eq:normal}
  \sqrt n(\hat F-F_0)
  \dto
  N(0,\Omega_F),
  \qquad
  \Omega_F=\Var(2e'Wz+e'W_1e).
\end{equation}
The $n^{-1}$ bias correction improves centering but does not change the
root-$n$ variance. By the delta method,
\[
  \sqrt n(\hat\varepsilon-\varepsilon_0)
  \dto
  N\!\left(0,\frac{\Omega_F}{4df^2\varepsilon_0^2}\right),
  \qquad F_0>0 .
\]
Thus the noncentral-$\chi^2$ interval belongs to the exact-fit/local-drift
regime. Under fixed misspecification, a Wald interval based on
\eqref{eq:normal} is the natural large-sample object.

On the point-estimation scale, write
\[
  c_{\mathrm{fix}}
  =
  \tr(\widetilde U\Gamma)+\beta_W+2e'Wb .
\]
Then a fixed-misspecification bias-corrected RMSEA estimator is
\begin{equation}\label{eq:rmsea}
  \hat\varepsilon_{\mathrm{fix}}
  =
  \sqrt{
    \max\!\left(\frac{\hat F-c_{\mathrm{fix}}/n}{df},0\right)
  } .
\end{equation}
Equivalently, if $T=n\hat F$, subtract $c_{\mathrm{fix}}$ from $T$, not
automatically $df$ or $\tr(U\Gamma)$.

\begin{remark}
The truncation in \eqref{eq:rmsea} is asymptotically irrelevant when $F_0>0$.
It only protects finite samples and connects the formula back to the usual
RMSEA convention near exact fit, where the limiting law changes from the normal
law in \eqref{eq:normal} to a quadratic form.
\end{remark}

\section{Estimator Cases}

\paragraph{ULS or fixed-weight GLS.}
This is the clean case. The weight is fixed, so $\beta_W=0$. If the saturated
moments are centered so that $b=0$, the correction is
$\tr(\widetilde U\Gamma)$. This is the most direct target for simulation checks:
compare $n\{\E[\hat F]-F_0\}$ with both $\tr(U\Gamma)$ and
$\tr(\widetilde U\Gamma)$.

\paragraph{Data-weighted GLS.}
The formula applies directly if $\hat W$ is a smooth function of $\hat m$ and
does not also depend on $\theta$. The new object is $\beta_W$, which is at least
linear in the residual through $g=W_1e$, with quadratic pieces through
$g'Pg$ and $e'W_2e$. It vanishes at exact fit but is part of the same $n^{-1}$
bias under fixed misspecification.

\paragraph{ML.}
Covariance ML is not simply a data-weighted quadratic form with a fixed
$\theta$-free weight. It can be locally represented by a quadratic form, but the
weight is tied to the model-implied covariance at the fitted parameter. That
means the same shape should reappear after a direct profile-likelihood
expansion, but the GLS $\beta_W$ term should not be copied over without doing
that calculation. The safe next derivation is therefore ML-specific, not a
blanket substitution into \eqref{eq:betaW}.

\section{What This Says}

The curvature correction is mathematically real. The fixed-misspecification
bias uses the profile Hessian $B=A-K$, so the trace is
$\tr(\widetilde U\Gamma)$ rather than $\tr(U\Gamma)$ unless either the residual
or the model curvature is negligible. This is the same phenomenon as replacing
expected information by observed Hessian in misspecification-robust standard
errors.

The $O(n^{-1})$ moment-bias term should be explicit. It is absent only after a
choice of first-stage convention, not as a theorem. For example, an unbiased
sample covariance convention removes it, while a different divisor may not.

The most useful next step is a small, fast validation harness for the fixed ULS
case. Once $\tr(\widetilde U\Gamma)$ is verified there, add controlled
data-weighted GLS examples to isolate $\beta_W$, and derive ML separately from
the profiled likelihood. That sequencing keeps the algebra testable: first
curvature, then weight influence, then likelihood-specific cancellation.

\end{document}
